\documentclass{article}
\usepackage{tikz}
\usetikzlibrary{angles, quotes}
\usepackage{amssymb}
\usepackage{pgfplots}
\pgfplotsset{width=10cm,compat=1.9}
\usepgfplotslibrary{statistics}
\usepackage[left=2.54cm,right=2.54cm,top=2.54cm,bottom=2.54cm]{geometry}
\usepackage{tasks} %Paquete necesario para la producción de items horizontales para algunos ejercicios de tarea empleando el comando /begin{multicols}}{#}.../end{multicols} para ayudarnos a reducir las páginas
\usepackage{pdflscape} %comando que permite cambiar a orientación horizontal las hojas%
\usepackage{soul} %Paquete para permitir la compilación de acentos usando el comando {\hl{}}}\\
\sethlcolor{yellow(munsell)} %comando para definir el color para el subrayado usando el comando \hl
\usepackage[utf8]{inputenc}
\usepackage[spanish]{babel}
\usepackage{icomma}
\usepackage{ragged2e}
\usepackage{siunitx}
\usepackage{url}
\usepackage[colorlinks=true, urlcolor=blue,  linkcolor=black, citecolor=green]{hyperref}
\usepackage{pdfpages}
\usepackage{blindtext} %Paquete que genera texto ficticio (dummy text) con el comando: \blindtext
\usepackage{cancel} %in the preamble gives you four different modes of striking through: \cancel{text to cancel} draws a diagonal line (slash) through its argument, \bcancel{text to cancel} uses the negative slope (a backslash), \xcancel{text to cancel} draws an X (actually \cancel plus \bcancel), \cancelto{〈value〉}{〈expression〉} draws a diagonal arrow through the 〈expression〉pointing to the 〈value〉 (math-mode only)
\usepackage{csquotes}
\usepackage{afterpage}
\usepackage{parskip} 
\usepackage{float}
\usepackage{enumitem}
\usepackage{multicol}%Paquete que permite la creación aislada de columnas de texto. De esta forma se reduce la cantidad de páginas en nuestro documento
\usepackage{lipsum} %Paquete que genera texto de relleno al igual que \usepackage{blindtext}, se usa con el comando \lipsum[#-#] los #(hashtags) delimitan el número de páginas que deseamos ocupar del paquete lipsum para usarlos como relleno. 

% Definir el comando personalizado para la nota
\newcommand{\mynote}[1]{%
\begin{tikzpicture}[baseline=-0.75ex]
\node[draw, circle, fill=Apple Green!60, inner sep=2pt] (note) {\textbf{Nota:}};
\end{tikzpicture}%
\ \textit{#1}%
}

\newenvironment{Figura}
{\par\medskip\noindent\minipage{\linewidth}}
{\endminipage\par\medskip}
\usepackage{caption}
\usepackage[
backend=biber,
sorting=none,
url=true
]{biblatex} %bibliografía
\addbibresource{biblio.bib}
% Establecer el espacio entre las entradas de la bibliografía
\setlength{\bibitemsep}{1\baselineskip} % Puedes ajustar el valor según tus preferencias
\usepackage{amsmath,amsthm,amssymb,amsfonts}
\usepackage{pifont} %Permite la compilación del comando \xmark: es el símbolo de la tachita
\usepackage{mathtools} %permite la compilación de símbolos para matrices
\usepackage{empheq} %Paquete que se relaciona con \usepackage[most]{tcolorbox} para la creación de cuadros/cajas de colores para delimitar los resultados matemáticos o líneas de texto usando la línea de comando: \begin{empheq}[box={\mymath[colback=orange(webcolor)!70,drop lifted shadow]}]{equation*} \end{empheq}
\usepackage[most]{tcolorbox} %Paquete que se relaciona con \usepackage{empheq} para la creación de cuadros/cajas de colores para delimitar los resultados matemáticos o líneas de texto
\renewcommand{\qedsymbol}{$\blacksquare$}
\usetikzlibrary{positioning,decorations.pathreplacing} %Paquete que permite la compilación de llaves para cuadros sinópticos (brace diagram) 
\usepackage{schemata} %The schemata package is designed just for these brace diagrams
\usepackage{booktabs}

\newcommand\AB[2]{\schema{\schemabox{#1}}{\schemabox{#2}}} %comando o paquete necesario para crear cuadros sinópticos

\newtcbox{\mymath}[1][]{%
nobeforeafter, math upper, tcbox raise base,
enhanced, colframe=blue!30!black,
colback=blue!30, boxrule=1pt,
#1} %Comando importante para encasillar los resultados matemáticos en cajas de diferentes colores, se relaciona con los paquetes \usepackage{empheq} y \usepackage[most]{tcolorbox} 

\newenvironment{sysmatrix}[1]
{\left(\begin{array}{@{}#1@{}}}
{\end{array}\right)}
\newcommand{\ro}[1]{%
\xrightarrow{\mathmakebox[\rowidth]{#1}}%
}
\newlength{\rowidth}% row operation width
\AtBeginDocument{\setlength{\rowidth}{3em}}  %Comando importante para laproducción de lineas de operaciones entre matrices, método de Gauss o Gauss Jordan se relaciona con \newenvironment{sysmatrix}

%formato para cambiar el horario a español
\usepackage[spanish]{datetime2}
\DTMsetdatestyle{spanish}

\renewcommand{\today}{\DTMdisplaydate{\the\year}{\the\month}{\the\day }{-1}}
%%%%%%%%%%%%%%%%%%%%%%%%%%%%%%%%%%%%%%%%%

\usepackage{datetime}
\newcommand{\mycurrenttime}{\xxivtime}

%%%%%%%%%%%%%%%%%%%%%%%%%%%%%%%%%%%%%%%%%%%%%%%%%%%%%%%%%%%%%%%%%%%%%%%%%%%%%%%%%%%%%%%%%%%%%%%%%%%%%%%%%%%%%%%%%%%%%%%%%%%%%%%%%%%%%%%%%%%%%%%%%


%%%%%%%%%%%%%Caja de color para el título%%%%%%%%%%%%%%%%%%%%%%%%%%%%%%%%%%%%%%%%%%%%%%%%%%%%%%

\definecolor{myframecolor}{RGB}{1, 45, 75} % Prussian Blue
\definecolor{myboxcolor}{RGB}{0, 107, 92} % Apple Green

\newtcolorbox{mybox}{
enhanced,
colback=myboxcolor!7, % Color de fondo del cuadro
colframe=myframecolor, % Color del marco del cuadro
arc=0pt,
boxrule=1pt,
borderline west={2mm}{-10mm}{myframecolor}, % Borde en el lado izquierdo
sharp corners=southwest,
width=\linewidth,
before=\par\vspace{\bigskipamount}, % Espacio antes del cuadro
after=\par\vspace{\bigskipamount} % Espacio después del cuadro
}


\newcommand{\euler}{e} %Comando para producir letra e de euler
\newenvironment{remark}{\par\vfill\footnotesize % Vertical white space above the remark and smaller font size
\begin{list}{}{
	\leftmargin=80pt % Indentation on the left
	\rightmargin=60pt}\item\ignorespaces % Indentation on the right
\makebox[-2.5pt]{\begin{tikzpicture}[overlay]
		\node[draw=Horizon!60,line width=2.5pt,circle,fill=Horizon!25,font=\sffamily\bfseries,inner sep=4pt,outer sep=0pt] at (-19pt,5pt){\textcolor{Horizon}{Nota}};\end{tikzpicture}} % Blue Nota in a circle
\advance\baselineskip -1pt}{\end{list}\vskip5pt} % Tighter line spacing and white space after remark
\usepackage{graphicx}

\usepackage{titling}

\input{colores.tex}

%%Producir signos de cita textual``Asignatura''

\newcommand{\xmark}{\ding{55}} %%%comando que genera la tachita

\newenvironment{MyColorPar}[1]{%
\leavevmode\color{#1}\ignorespaces%
}{%
}%

%%%%%%%%%%%%%%%%%%% Título de la tarea, Nombre de alumno

\title{Clase \# 37 - Mecánica Estadística}
\author{\emph{Notas de clase} $\boldsymbol{\mid}$}
\date{\today}

%%%%%%%%%%%%%%%%%%% Datos de la Materia

\usepackage{fancyhdr}
\fancypagestyle{plain}{%  the preset of fancyhdr 
\fancyhf{} % clear all header and footer fields
\fancyfoot[R]{
\includegraphics[width=2cm]{LogoFCFMBUAP (1).png}}
\fancyfoot[L]{{\bfseries{\thedate{}}} a las {\bfseries{\mycurrenttime{}}} horas{} (GMT-6, H. Puebla de Zaragoza, Pue)}
\fancyhead[L]{Mecánica Estadística}
\fancyhead[R]{\theauthor}
}
\makeatletter
\def\@maketitle{%
\newpage
\null
\vskip 1em%
\begin{center}%
\let \footnote \thanks
{\LARGE \@title \par}%
\vskip 1em%
%{\large \@date}%
\end{center}%
\par
\vskip 1em}
\makeatother

\usepackage{lipsum}  
\usepackage{cmbright}


%%%%%%%%%%%%%%%%%%%%%%%%%%%%%%%%%%%%%%%%%%%%%%%%%%%%%%%%%%%%%%%%%%%%%%%%%%%%%%%%%%%%%%%%%%%%%%%%%%%%%%%%%%%%%%%%%%%%%%%%%%%%%%%%%%%%%%%%%%%%%%%%%%%%%%%%%%%%%%%%%%%%%%%%%%%%%%%%%%%%%%%%%%%%%%%%%%%%%%

\begin{document}



\maketitle


%%%%%%%%%%%%%%%%%%% Datos del profesor, campus y aula

\begin{mybox}
\noindent\begin{tabular}{@{}ll}
	{\bfseries{Profesor}} & Dr. Jorge Velázquez Castro\\
	\textcolor{prussianblue}{\bfseries{Campus}} / \textcolor{trueblue}{\bfseries{Facultad}}  & \textcolor{prussianblue}{\bfseries{C.U. BUAP}} /  \textcolor{trueblue}{\bfseries{FCFM}} \\
	\textcolor{prussianblue}{\bfseries{Edificio}} /  \textcolor{trueblue}{\bfseries{Salón}}    & 	\textcolor{prussianblue}{\bfseries{EMA4}} /  \textcolor{trueblue}{\bfseries{407}}   
\end{tabular} 
\end{mybox} \vspace{1cm}

\tableofcontents

\vspace{2cm}

\section{Clase de hoy: Gases con grados de libertad internos}

\begin{figure}[H]
    \centering
    \includegraphics[width=0.9\textwidth]{1.png}
    \caption{Esquema general de un gas con grados de libertad internos.}
    \label{fig:derivada_total_H}
\end{figure}

La \textbf{Hamiltoniana} del sistema anterior corresponde a la suma de todas las partículas, es decir, una suma sobre la variable $n$:

\[
\mathcal{H} = \sum\limits_{n} \left[ \mathcal{H}_{\text{tras}}^{(n)} + \mathcal{H}_{\text{rot}}^{(n)} + \mathcal{H}_{\text{vib}}^{(n)} \right]
\]

Como el Hamiltoniano está definido como una suma, en el \textbf{ensamble canónico} podemos aplicar la propiedad de \textbf{factorización}. De este modo, la función de partición total se expresa como:

\[
Z(T,V,N) = \frac{1}{N!}\, Z_{\text{tras}}^{N}\, Z_{\text{rot}}^{N}\, Z_{\text{vib}}^{N}
\]

donde cada $Z$ representa la función de partición correspondiente (traslacional, rotacional o vibracional).  
Por lo tanto, la \textbf{energía libre de Helmholtz} puede escribirse como:

\[
F(T,V,N) = F_{\text{tras}}(T,V,N) + F_{\text{rot}}(T,V,N) + F_{\text{vib}}(T,V,N)
\]

\subsection{Parte traslacional}

El Hamiltoniano traslacional es:

\[
\mathcal{H}_{\text{tras}} = \frac{\vec{p}^{\,2}}{2m}
\]

y su función de partición:

\[
Z_{\text{tras}} = V \left( \frac{2\pi m k_{B} T}{h^{2}} \right)^{\!\frac{3}{2}}
\]


\subsection{Parte rotacional}

\begin{figure}[H]
    \centering
    \includegraphics[width=0.6\textwidth]{2.png}
    \caption{Cuerpo rígido con momento de inercia $I$.}
    \label{fig:rotacional}
\end{figure}

Existe un momento de inercia asociado a la rotación, una característica propia del cuerpo rígido.  
Para no complicar el análisis, asumimos que la molécula posee un solo momento de inercia relevante, mientras que los otros son pequeños y pueden despreciarse. En otras palabras, consideramos que sólo rota alrededor de un eje principal.

El Hamiltoniano rotacional se escribe como:

\[
\mathcal{\hat{H}}_{rot} = \frac{\hat{L}^{2}}{2I}
\]

donde el acento indica que puede tratarse tanto del caso clásico como del cuántico.  
$I$ representa el \textcolor{Apple Green}{\textbf{momento de inercia}} y $L$ el \textcolor{Sun}{\textbf{momento angular}} .

Cuando analizamos sistemas cuánticos, las variables clásicas se promueven a \textbf{operadores}, y los valores medibles corresponden a los \textcolor{Cinnabar}{\textbf{autovalores}} de dichos operadores, determinados por las condiciones de frontera del sistema.

Así, los \textbf{niveles energéticos} resultan ser:

\[
\epsilon_{\ell} = \frac{\ell (\ell +1)\hbar^{2}}{2I}, 
\qquad \ell = 0,1,2,\ldots
\]

Además, aparece el concepto de \textbf{degeneración}.  
Podemos imaginar un cuerpo rotando: existen distintas maneras de hacerlo con la misma energía.  
La degeneración correspondiente es:

\[
\omega_{\ell} = 2\ell +1
\]

Esto se interpreta como el número de orientaciones posibles de rotación para cada $\ell$.  
Por ejemplo, a energías bajas ($\ell=1$) hay tres formas de giro: hacia arriba, horizontal y hacia abajo.  
A medida que el momento angular aumenta, crece también el número de formas posibles, hasta que el comportamiento se asemeja al de un cuerpo macroscópico.

Por tanto, la función de partición rotacional es:

\[
Z_{\text{rot}} = \sum\limits_{\ell=0}^{\infty} (2\ell +1)\,
e^{ -\beta \frac{\ell (\ell +1)\hbar^{2} }{2I} }
\]

Esta suma no puede resolverse de \underline{manera cerrada}, por lo que se aproxima mediante una integral, válida en el régimen de \textbf{altas temperaturas}.


\subsubsection*{Aproximación de Euler--Maclaurin}

Para altas temperaturas utilizamos la fórmula de \textbf{Euler--Maclaurin}:

\[
\sum\limits_{\ell =0}^{\infty} f(\ell) \approx
\int_{0}^{\infty} f(x)\,dx + \frac{1}{2}f(0)
- \frac{1}{12} f^{\prime}(0) + \frac{1}{720} f^{\prime\prime\prime}(0) + \ldots
\]

donde:
\[
f(x) = (2x+1)e^{ -\frac{\Theta}{T}x(x+1) },
\qquad
\Theta = \frac{\hbar^{2}}{2I k_{B}}
\]
es la \textbf{temperatura característica} del sistema.

Evaluando la integral principal:

\[
\int_{0}^{\infty} f(x)\,dx
= \int_{0}^{\infty} (2x+1)e^{- \frac{\Theta}{T} x(x+1) } dx
\simeq \frac{T}{\Theta}
\]

Aplicando los términos de corrección:

\[
Z_{\text{rot}} \approx 
\frac{T}{\Theta} + \frac{1}{3} + \frac{1}{15}\frac{\Theta}{T}
+ \frac{4}{315}\left( \frac{\Theta}{T} \right)^{2} + \ldots
\]

A \textbf{temperaturas altas} ($T\to\infty$):

\[
\boxed{Z_{\text{rot}} \approx \frac{T}{\Theta}
= \frac{2I k_{B}T}{\hbar^{2}}}
\]

Mientras que a \textbf{temperaturas bajas} sólo se conservan los primeros términos:

\[
Z_{\text{rot}}(T\to 0) \approx
1 + 3e^{ - \frac{2\Theta}{T} } + 5e^{ - \frac{6\Theta}{T} } + \ldots
\]

\subsection{Parte vibracional}

\begin{figure}[H]
    \centering
    \includegraphics[width=0.9\textwidth]{4.png}
    \caption{Modos vibracionales de una molécula.}
    \label{fig:vibracional}
\end{figure}

La \textcolor{vividviolet}{\textbf{Hamiltoniana vibracional}}  es:

\[
\mathcal{H}_{\text{vib}} = 
\sum\limits_{i=1}^{f} \left(
\frac{p_{i}^{2}}{2B_{i}} + 
\frac{1}{2}B_{i}\omega^{2}q_{i}^{2}
\right)
\]

donde $f$ son los modos de vibración de la molécula, $B_{i}$ la masa reducida y $\omega$ la frecuencia del modo.

La función de partición vibracional resulta:

\[
Z_{\text{vib}} = 
\prod\limits_{i=1}^{f} 
\left[ 
2\sinh\!\left( 
\frac{1}{2}\beta\hbar\omega_{i}
\right)
\right]^{-1}
\]

y el calor específico asociado a los modos vibracionales es:

\[
C_{V}^{(\text{vib})} = 
\sum\limits_{i=1}^{f}k_{B}(\beta\hbar\omega_{i})^{2}
\frac{e^{\beta\hbar\omega_{i}}}{( e^{\beta\hbar\omega_{i}} -1 )^{2}}
\]

\subsection{Interpretación física y discusión (\textcolor{orange(colorwheel)}{\textbf{¡importante!}}) }

\begin{figure}[H]
    \centering
    \includegraphics[width=1.1\textwidth]{5.png}
    \caption{Comportamiento del calor específico con la temperatura.}
    \label{fig:cv}
\end{figure}

En los experimentos, la forma de la curva del calor específico $\boldsymbol{C_{V}}$ muestra cómo distintos grados de libertad se van activando con la temperatura.

A bajas temperaturas, sólo se aprecian los grados de libertad \textbf{traslacionales}, con un $C_{V} = \tfrac{3}{2}k_{B}$.  
Al aumentar la temperatura, comienzan a excitarse los grados de libertad \textbf{rotacionales}, lo que añade aproximadamente $k_{B}$ adicional, dando $C_{V} = \tfrac{5}{2}k_{B}$ en gases diatómicos.

En moléculas \textbf{triatómicas} se suman más contribuciones, pues existen más ejes de rotación y modos vibracionales.  
Cada nuevo grado de libertad aporta aproximadamente otra cantidad de orden $k_{B}T$, según el teorema de la equipartición.

Los grados de libertad \textbf{vibracionales} no se observan a bajas temperaturas porque la mayoría de las moléculas permanecen en el nivel vibracional más bajo.  
Al aumentar la temperatura, estos niveles comienzan a excitarse, y eventualmente también los grados de libertad \textbf{electrónicos}.  
A energías aún mayores podrían activarse los grados de libertad \textbf{nucleares}, e incluso los asociados a las \textbf{fuerzas fuertes} dentro de los quarks.

Por eso, los \textbf{aceleradores de partículas} pueden interpretarse metafóricamente como verdaderos “\underline{hornos}” que permiten alcanzar energías tan altas que revelan los grados de libertad internos más profundos de la materia.

Como un ejercicio de práctica (\underline{no se califica	}) el profesor sugiere que probemos valores para graficar y ver el comportamiento, puede ser en \textcolor{Sun}{\textbf{python}}, \textcolor{Logan}{\textbf{matlab}}, \textcolor{Tarawera}{\textbf{Wolfram}} etc. 

\subsection{Aviso como recordatorio ya hay tarea en la plataforma. Examen 7 de noviembre de lo de ensemble canónico.}


\end{document}


